\documentclass[10pt,a4paper]{amsart}
\usepackage[margin=19mm]{geometry}
\usepackage{amsmath,amssymb,mathtools}
\usepackage[scr=rsfs,cal=euler]{mathalfa}
\usepackage{xfrac}
\usepackage[unicode,hidelinks,bookmarks=false]{hyperref}
\newcommand{\ie}{i.e.\ }
\newcommand{\cf}{cf.\ }
\newcommand{\resp}{respectively\ }
\newcommand{\Subsection}{\subsection}
\providecommand{\nobreakdash}{\nobreak\hbox{-}}
\setlength{\emergencystretch}{2em}
\multlinegap0pt
\usepackage{amssymb}
\newtheorem{lemma}{Lemma}
\title{Gaussian consensus processes and their Lyapunov exponents}
\author{Edward Crane, Stanislav Volkov}
\date{Source: arXiv:2405.04951v2 (CC BY 4.0)}
\begin{document}
\maketitle

\noindent\emph{Source and licence.} Gaussian consensus processes and their Lyapunov exponents. Version arXiv:2405.04951v2 (CC BY 4.0), distributed by arXiv under the Creative Commons Attribution 4.0 International licence, http://creativecommons.org/licenses/by/4.0/, as recorded in the arXiv licence metadata for this exact version and shown on the abstract page https://arxiv.org/abs/2405.04951v2. Full licence notice: \texttt{licenses/arxiv-2405.04951-CC-BY-4.0.txt}.\par\smallskip

\noindent\textit{Editorial scope.} Counted unit: the article's lemma lem:phi (source0.tex lines 1462--1470). The article's complete original proof of it is delivered with it (source0.tex lines 1471--1536, one \texttt{proof} environment of 66 lines). No mathematics is written, repaired or re-derived: the statement and the proof are the article's own lines, in the article's own notation, and no retained line is substituted or re-flowed.  Retained uncounted as the source-local context the proof needs: lines 149--151.  Nothing was omitted from any retained span, so the package carries no omission record.  No retained line contains a citation, so the wrapper carries no bibliography and no bibliography entry.  No retained line contains a figure environment, an image-inclusion command or drawing code, so no image is delivered and the package declares no asset dependency.  Editorial formatting only, in addition: an AMS \texttt{amsart} preamble that restates the article's own declarations and macros used by the retained lines, namely the environments \texttt{lemma} on separate counters, together with the standard packages those lines need.  The retained environments print in this document as Lemma 1 in order of appearance, whereas the article prints the same items with the numbering of its own full document.  The complete native source of the article as distributed in the arXiv e-print for this exact version is bundled unchanged as \texttt{sources/158842/source0.tex}.
\begin{equation} \label{def of phi}
\phi_m(x) := e^{-x} \sum_{j=0}^\infty \frac{x^j}{j!} \psi(j+m)=\psi(m)+\int_0^1 \frac{(1-e^{-xs})(1-s)^{m-1}}{s}\mathrm{d}s\,,
\end{equation}

\begin{lemma}\label{lem:phi}
For all $m>0$ we have
\begin{align*}
\phi_m(x) &=\psi(m)
-\sum_{k=1}^\infty \frac{(-x)^k}{k\cdot m(m+1)\dots(m+k-1)}
=
\psi(m)+\int_0^1 \frac{(1-e^{-xs})(1-s)^{m-1}}{s}\mathrm{d}s.
\end{align*}
\end{lemma}

\begin{proof}
Differentiating the expression in~\eqref{def of phi}, and using the fact that $\psi(y+1)-\psi(y)=1/y$, we get
\begin{align}\label{eq:phiprime}
\begin{split}
\phi_m'(x)&=-e^{-x}\psi(m)-
e^{-x}\sum_{j=1}^\infty\psi(j+m)\left[\frac{x^j}{j!}-\frac{x^{j-1}}{(j-1)!}\right]
\\ &=
e^{-x}\left[ -\psi(m)+\psi(m+1)\right]
+e^{-x}\sum_{j=1}^\infty\frac{x^j}{j!}\left[\psi(m+j+1)-\psi(m+j)\right]
\\ &
=e^{-x}\left[\frac1m
+\sum_{j=1}^\infty\frac{x^j}{j!(m+j)}
\right]
=e^{-x}\sum_{j=0}^\infty\frac{x^j}{j!(m+j)}
=\mathbb{E}\left(\frac1{m+\xi}\right)
\end{split}\end{align}
where $\xi \sim \mathrm{Poi}(x)$. Note also that
$$
\left(x^m e^x \phi_m'(x)\right)'=\sum_{j=0}^\infty \frac{x^{j+m-1}}{j!}=x^{m-1} e^x
$$
so that $\phi_m(x)$ satisfies the second-order linear differential equation
$$
x \phi_m''+(x+m)\phi_m'=1;\quad\phi_m(0)=\psi(m); \quad \phi_m'(0)=\frac1m.
$$
The expression for $\phi_m'(x)$ can be simplified
further:
\begin{align*}
\phi_m'(x)&=\frac{e^{-x}}m\left[1
+\sum_{j=1}^\infty\frac{x^j}{(j-1)!} \left(\frac1j-\frac1{j+m}\right)
\right]=
\frac{e^{-x}}m\left[\sum_{j=0}^\infty\frac{x^j}{j!} 
-x\sum_{j=1}^\infty\frac{x^{j-1}}{(j-1)!(j+m)} 
\right]
\\ &=
\frac{e^{-x}}m\left[e^x
-x\sum_{j=0}^\infty\frac{x^{j}}{j!(j+m+1)} 
\right]
=\frac1m -\frac{x}m\phi'_{m+1}(x).
\end{align*}
Iterating this formula $\ell$ times, we obtain that
$$
\phi_m'(x)=\frac1m-\frac{x}{m(m+1)}
+\frac{x^2}{m(m+1)(m+2)}-\dots  \pm \frac{x^\ell}{m(m+1)\dots(m+\ell-1)}\phi'_{m+\ell}(x).
$$
Noting that $0\le \phi_m'(x)\le 1$ for all $m$ and letting $\ell\to\infty$ we obtain 
\begin{align}\label{phiT}
\phi_m'(x)=
\sum_{j=0}^\infty \frac{(-x)^j}{m(m+1)\dots(m+j)}.
\end{align}
Recalling that $\phi_m(0)=\psi(m)$, we obtain $$\phi_m(x)=\psi(m)-\sum_{j=0}^\infty \frac{(-x)^{j+1}}{(j+1)\cdot m(m+1)\dots(m+j)}.$$
To get the second expression for $\phi_m(x)$, note that
$$
\frac{1-e^{-xs}}{s}=x-\frac{x^2s}{2!}+\frac{x^3s^2}{3!}-\dots=-\sum_{k=1}^\infty \frac{(-x)^ks^{k-1}}{k!}
$$
and thus
\begin{align*}
\int_0^1 \frac{(1-e^{-xs})(1-s)^{m-1}}{s}\mathrm{d}s
&=-\int_0^1 \left(\sum_{k=1}^\infty \frac{(-x)^ks^{k-1}}{k!}\right)(1-s)^{m-1}\mathrm{d}s
\\ &
=-\sum_{k=1}^\infty \frac{(-x)^k}{k\,\Gamma(k)}\int_0^1 s^{k-1}(1-s)^{m-1}\mathrm{d}s
\\ &=-
\sum_{k=1}^\infty \frac{(-x)^k}{k\,\Gamma(k)} \frac{\Gamma(k)\,\Gamma(m)}{\Gamma(m+k)}
=-
\sum_{k=1}^\infty \frac{(-x)^k}{k\cdot m(m+1)\dots(m+k-1)}.
\end{align*}
\end{proof}

\end{document}
